% ─────────────────────────────────────────────────────────────────────────────
%  Chapter 5 — Remote Cluster Deployment
% ─────────────────────────────────────────────────────────────────────────────

% TikZ styles for this chapter's flowcharts
% NOTE: \gls{} is intentionally NOT used inside TikZ node labels — hyperref
%       link anchors conflict with pgf boxes and \scalebox.
\tikzset{
  rmproc/.style={rectangle, rounded corners=3pt, draw=black!70, fill=blue!8,
                 text width=5.8cm, minimum height=0.75cm, text centered,
                 font=\small, inner sep=5pt},
  rmdec/.style={diamond, aspect=2.4, draw=black!70, fill=orange!15,
                text width=3.0cm, minimum height=0.75cm, text centered,
                font=\small, inner sep=2pt},
  rmio/.style={ellipse, draw=black!70, fill=green!20,
               text width=2.6cm, minimum height=0.7cm,
               text centered, font=\small, inner sep=3pt},
  rmside/.style={rectangle, rounded corners=3pt, draw=black!60, fill=gray!10,
                 text width=4.0cm, minimum height=0.75cm, text centered,
                 font=\small, inner sep=5pt},
  rmblock/.style={rectangle, rounded corners=4pt, draw=black!70,
                  text width=3.6cm, minimum height=1.1cm, text centered,
                  font=\small, inner sep=6pt},
  rmclient/.style={rectangle, rounded corners=4pt, draw=black!70, fill=cyan!8,
                   text width=4.5cm, minimum height=1.4cm, text centered,
                   font=\small, inner sep=6pt},
  rmvm/.style={rectangle, rounded corners=4pt, draw=black!70, fill=orange!10,
               text width=4.5cm, minimum height=1.4cm, text centered,
               font=\small, inner sep=6pt},
  rmarr/.style={thick,->,>=Stealth},
  rmssh/.style={thick,->,>=Stealth,dashed,blue!60},
  rmsftp/.style={thick,->,>=Stealth,dotted,green!60!black},
  rmline/.style={thick},
}

Chapter~4 validated the \gls{aerosouda} concept through a local cluster prototype running on a
single workstation.  That prototype demonstrated that the proposed distributed MapReduce pipeline
is technically feasible.  The natural next step is to extend the deployment model so that the
Spark and \gls{hdfs} cluster runs on a dedicated remote server with more computational resources,
while the analyst interacts with it from a lightweight client on their own machine without
requiring Docker to be installed locally.

This chapter describes the remote cluster deployment architecture, implemented in the
\texttt{tg-cluster} repository.  The repository is split into two independent components:
\texttt{cluster/}, the Docker Compose stack that runs on the remote \gls{vm},
and \texttt{client/}, a Python-only \gls{cli} that runs on the analyst's machine and communicates
with the \gls{vm} exclusively via \gls{ssh} and \gls{sftp}.

The chapter is organized as follows.
Section~\ref{sec:rm-overview} presents the two-component architecture and justifies the
design decisions.
Section~\ref{sec:rm-cluster} documents the \gls{vm} cluster configuration.
Section~\ref{sec:rm-client} describes the remote client \gls{cli} and its eight-step pipeline.
Section~\ref{sec:rm-ssh-sftp} explains the \gls{ssh}/\gls{sftp} communication design.
Section~\ref{sec:rm-dataflow} presents the end-to-end data flow from the analyst's machine to
the \gls{vm} cluster and back.

% ─────────────────────────────────────────────────────────────────────────────
\section{Architecture Overview}
\label{sec:rm-overview}
% ─────────────────────────────────────────────────────────────────────────────

\subsection{Motivation for a Two-Component Design}
\label{subsec:rm-motivation-two-component}

The local cluster prototype from Chapter~4 requires the analyst to have Docker installed and to
run the entire Spark and \gls{hdfs} stack on their own machine.  This is acceptable for
development and testing but impractical for production use: simulation result \gls{csv} files
can reach several gigabytes; a local workstation may lack the memory and processing cores to
distribute such heavy workload effectively; and the cluster must be shut down and restarted
each time the analyst switches computers.

A remote cluster hosted on a \gls{vm} solves all three problems: the \gls{vm} can be provisioned
with many cores and ample memory, it stores the simulation \gls{csv} files persistently on
\gls{hdfs}, and it is accessible from any machine with an internet connection.  The analyst only
needs Python and a handful of pure-Python dependencies on their laptop in order to run this
repository. This project also includes in its \texttt{README} file a step-by-step guide to
install said python dependencies.

\subsection{Repository Structure}
\label{subsec:rm-repository-structure}

The \texttt{tg-cluster} repository is organized into two top-level directories:

\begin{itemize}
  \item \texttt{cluster/}, the Docker Compose stack intended to be deployed on the remote
        \gls{vm}.  It contains the same Spark and \gls{hdfs} containers as the local cluster,
        together with the processing scripts, scenario \gls{yaml} files, and all configuration
        files.  Analysts clone this directory on the \gls{vm} and start it with
        \texttt{make~up}.
  \item \texttt{client/}, the analyst-side \gls{cli}.  It has no Docker dependency; its only
        requirements are Python~3.8+ and the packages listed in \texttt{requirements.txt} (rich,
        questionary, pyyaml, pandas, paramiko).  The analyst runs \texttt{python~cli.py} on their
        laptop to initiate the full distributed pipeline on the remote \gls{vm}.
\end{itemize}

Figure~\ref{fig:rm-arch} illustrates the two-component architecture and the communication
channels between them.

\begin{figure}[htbp]
\centering
\begin{tikzpicture}[x=1cm, y=1cm]

%% Client box
\node[rmclient] (client) at (-4.5, 0)
    {Analyst machine\\\texttt{client/cli.py}\\(Python only, no Docker)};

%% VM box
\node[rmvm] (vm) at (4.5, 0)
    {Remote VM\\\texttt{cluster/} (Docker)\\Spark + HDFS};

%% SSH arrow
\draw[rmssh] (client.east) to[bend left=18]
    node[above,font=\scriptsize]{SSH: job submission} (vm.west);

%% SFTP arrow
\draw[rmsftp] (vm.west) to[bend left=18]
    node[below,font=\scriptsize]{SFTP: result files} (client.east);

%% Client local files
\node[rmblock, fill=cyan!5] (lcfiles) at (-4.5, -2.6)
    {\texttt{scenarios/} (YAML)\\\texttt{outputs/} (results)};
\draw[rmarr] (client.south) -- (lcfiles.north);

%% VM HDFS
\node[rmblock, fill=orange!5] (hdfs) at (4.5, -2.6)
    {\texttt{HDFS /inputs/}\\\texttt{(simulation CSVs)}};
\draw[rmarr] (vm.south) -- (hdfs.north);

\end{tikzpicture}
\caption[Two-component architecture of the AeroSOuDA remote deployment.]{Two-component architecture of the AeroSOuDA remote deployment.
         Dashed blue arrows indicate SSH connections; dotted green arrows indicate SFTP
         file transfers.  YAML scenario files reside on the analyst's machine; simulation
         CSV files reside on the VM's HDFS.}
\label{fig:rm-arch}
\end{figure}

% ─────────────────────────────────────────────────────────────────────────────
\section{The VM Cluster (\texttt{cluster/})}
\label{sec:rm-cluster}
% ─────────────────────────────────────────────────────────────────────────────

\subsection{Container Configuration}
\label{subsec:rm-container-configuration}

The \texttt{cluster/docker-compose.yml} defines the same thirteen-container topology as the
local cluster: one \gls{hdfs} NameNode, four \gls{hdfs} DataNodes, one Spark master, six Spark
workers, and one Spark History Server.

Whereas the local cluster is started by the analyst on
their own workstation and torn down after each session, the \gls{vm} cluster is intended to run
continuously.  Simulation result \gls{csv} files are uploaded to the \gls{hdfs} storage layer
once and remain there for repeated analysis runs.

\subsection{Persistent HDFS Storage}
\label{subsec:rm-persistent-hdfs-storage}

One of the main advantages of the remote deployment is that large \gls{csv} files are uploaded
to the \gls{vm}'s \gls{hdfs} once and do not need to be retransferred for each analysis run.
The simulation result files can be copied to the \gls{vm} via the Secure Copy Protocol
(\texttt{scp}) (or a similar file transfer tool), then loaded into \gls{hdfs} with:

\begin{verbatim}
docker cp results.csv namenode:/tmp/results.csv
docker exec namenode hdfs dfs -put /tmp/results.csv /inputs/
\end{verbatim}

After this one-time upload, any number of analysis runs with different \gls{yaml} scenario
configurations can be performed without retransferring the \gls{csv}.

\subsection{VM Setup}
\label{subsec:rm-vm-setup}

Deploying the cluster on a fresh \gls{vm} requires only Docker and Docker Compose.  The setup
procedure is documented in \texttt{cluster/README.md} and consists of three steps:

\begin{enumerate}
  \item Install Docker Engine and the Compose plugin following the official Docker documentation.
  \item Clone the \texttt{tg-cluster} repository and navigate to \texttt{cluster/}.
  \item Run \texttt{make fix-perms \&\& make up} to fix volume permissions and start all
        containers.
\end{enumerate}

The \gls{hdfs} web interface is available at \texttt{http://<VM\_IP>:9870} and the Spark web
interface at \texttt{http://<VM\_IP>:8080}.  Both interfaces should be accessible from the
analyst's machine if the \gls{vm} firewall allows \gls{tcp} connections on those ports.

% ─────────────────────────────────────────────────────────────────────────────
\section{The Remote Client (\texttt{client/})}
\label{sec:rm-client}
% ─────────────────────────────────────────────────────────────────────────────

\subsection{Overview}
\label{subsec:rm-client-overview}

The remote client (\texttt{client/cli.py}) is the analyst-facing entry point for the
remote cluster deployment.  It provides the same interactive experience as the local \gls{cli}
but replaces all Docker commands with \gls{ssh} and \gls{sftp} operations directed at the remote
\gls{vm}.  The client has no Docker dependency; it uses the \texttt{paramiko} library
\parencite{paramiko2024} for \gls{ssh} connectivity and the standard \texttt{pyyaml} and
\texttt{rich} libraries for scenario management and terminal output.

\subsection{Pipeline Flow}
\label{subsec:rm-pipeline-flow}

Figure~\ref{fig:rm-flow} presents the remote client pipeline. The pipeline is composed of eight
sequential stages.

\begin{figure}[p]
\centering
\scalebox{0.78}{%
\begin{tikzpicture}[x=1cm, y=1cm]

\node[rmio]   (start) at (0,  0.0) {Start};
\node[rmproc] (wel)   at (0, -1.4) {1. Display welcome screen};
\node[rmproc] (conn)  at (0, -2.8) {2. Prompt VM connection (IP, user, auth)};
\node[rmproc] (ssh)   at (0, -4.2) {3. Connect via SSH (paramiko)};
\node[rmproc] (hdfs)  at (0, -5.6) {4. List CSVs from VM HDFS via SSH};
\node[rmproc] (pick)  at (0, -7.0) {5. Pick CSV + pick or create YAML scenario};
\node[rmproc] (conf)  at (0, -8.4) {6. Confirm execution};
\node[rmproc] (job)   at (0, -9.8) {7. Submit job via SSH; stream progress};
\node[rmproc] (sftp)  at (0,-11.2) {8. SFTP: download results to local outputs/};
\node[rmproc] (disp)  at (0,-12.6) {9. Display results};
\node[rmio]   (end)   at (0,-14.0) {End};

\draw[rmarr] (start) -- (wel);
\draw[rmarr] (wel)   -- (conn);
\draw[rmarr] (conn)  -- (ssh);
\draw[rmarr] (ssh)   -- (hdfs);
\draw[rmarr] (hdfs)  -- (pick);
\draw[rmarr] (pick)  -- (conf);
\draw[rmarr] (conf)  -- (job);
\draw[rmarr] (job)   -- (sftp);
\draw[rmarr] (sftp)  -- (disp);
\draw[rmarr] (disp)  -- (end);

\end{tikzpicture}%
}
\caption{Execution pipeline of the AeroSOuDA remote client (\texttt{client/cli.py}).}
\label{fig:rm-flow}
\end{figure}

The stages of the pipeline are described as follows:

\begin{enumerate}

\item Welcome screen. A styled panel introduces the tool, identifies the
      remote-de\-ploy\-ment context, and lists the pipeline steps.

\item VM connection prompt. The client checks for a \texttt{.env} file in the
      working directory containing the five connection variables: \texttt{SSH\_IP},
      \texttt{SSH\_USERNAME}, \texttt{SSH\_PASSWORD}, \texttt{SSH\_PORT}, and
      \texttt{VM\_CLUSTER\_PATH}.  If the file is present and all variables are
      set, all connection parameters are loaded automatically and no prompt is shown.
      Otherwise, the analyst is prompted interactively for the \gls{vm} \gls{ip} address,
      \gls{ssh} username, authentication method (password or private key file path), port,
      and cluster directory path.  The \texttt{.env} approach is the recommended workflow
      for repeated use against the same cluster, as it eliminates repeated manual input.

\item SSH connection. A \texttt{paramiko.SSHClient} is opened to the \gls{vm}.  The
      connection module then issues a \texttt{docker ps} command over \gls{ssh} to verify that
      the cluster containers (\texttt{spark-master}, \texttt{namenode}, DataNodes, and workers)
      are running.  If any required container is absent, the pipeline exits with actionable
      guidance.

\item CSV listing from HDFS. The client issues the command
      \texttt{docker exec namenode hdfs dfs -ls /inputs/} over \gls{ssh} and parses the output
      to build a list of available \gls{csv} files together with their sizes.  This listing
      reflects the current contents of the \gls{hdfs} storage layer on the \gls{vm}.

\item CSV and scenario selection. The analyst selects one of the listed \gls{csv} files
      and either loads a \gls{yaml} scenario from the local \texttt{scenarios/} directory or
      creates a new one using the same guided wizard as the local \gls{cli}.

\item Execution confirmation. A summary panel displays the selected \gls{csv} name,
      the \gls{vm} host, and the scenario configuration.  The analyst confirms before the job
      is submitted.

\item Job submission and monitoring. The client performs the following sub-steps:
  \begin{enumerate}[label=\alph*.]
    \item \gls{sftp}: upload the local \gls{yaml} scenario to \texttt{/tmp/} on the \gls{vm}.
    \item \gls{ssh}: inject the \gls{yaml} into the \texttt{spark-master} container via
          \texttt{docker exec}.
    \item \gls{ssh}: execute \texttt{run\_csv\_processor.sh} inside \texttt{spark-master},
          streaming \texttt{stdout} back through the \gls{ssh} channel.  Spark stage-progress
          lines are parsed and displayed as a live progress bar on the analyst's terminal.
    \item \gls{ssh}: copy output artefacts from the container to a temporary directory on the
          \gls{vm} via \texttt{docker cp}.
  \end{enumerate}

\item Result transfer. The client uses \gls{sftp} to download all output files from
      the \gls{vm}'s temporary directory to the local \texttt{outputs/} directory.  The three
      output artefacts an aggregated \gls{csv}, a scatter plot image, and a \gls{json}
      summary are saved locally on the analyst's machine.

\item Results display. A final panel reports the scenario name, the number of data
      points, the batch aggregate value, and the local paths of the output files.

\end{enumerate}

% ─────────────────────────────────────────────────────────────────────────────
\section{SSH and SFTP Communication Design}
\label{sec:rm-ssh-sftp}
% ─────────────────────────────────────────────────────────────────────────────

\subsection{The \texttt{RemoteCluster} Class}
\label{subsec:rm-remotecluster-class}

\begin{sloppypar}
All \gls{ssh} and \gls{sftp} operations are encapsulated in the
\texttt{cli/remote.py::\allowbreak{}RemoteCluster} class.  On \texttt{connect()}, it opens a single
\texttt{paramiko.SSHClient} connection to the \gls{vm} and derives an \texttt{SFTPClient}
from the same underlying transport, minimizing the number of \gls{tcp} handshakes.
\end{sloppypar}

For job execution, \texttt{run\_job()} opens a pseudo-terminal channel on the existing transport.
The pseudo-terminal is required because the \texttt{docker compose exec} command on the \gls{vm}
buffers output differently when it detects that its standard output is not a \gls{tty}.
Allocating a pseudo-terminal forces line-buffered output, which is necessary for the stage-progress
lines to appear in real time on the analyst's terminal.

\subsection{Authentication}
\label{subsec:rm-authentication}

The \texttt{RemoteCluster} constructor accepts either a password or the path to an \gls{ssh}
private key file.  When a key path is provided, \texttt{paramiko} loads the key with
\texttt{para\-miko.\allowbreak{}RSAKey.\allowbreak{}from\_private\_key\_file()}.
Both authentication methods are equivalent
in security from \texttt{paramiko}'s perspective; using a private key is recommended for
production deployments because it avoids storing or transmitting passwords.

\subsection{Data Transfer Design}
\label{subsec:rm-data-transfer-design}

Three distinct data-transfer patterns are used in the pipeline:

\begin{enumerate}
  \item YAML upload (\gls{sftp} PUT): the analyst's local \gls{yaml} scenario file is
        uploaded to \texttt{/tmp/aerosouda\_scenario.yaml} on the \gls{vm} filesystem, then
        injected into the \texttt{spark-master} container via a Python one-liner piped through
        \texttt{docker exec -i}.
  \item Job output export (\gls{ssh} command): after the Spark job completes, the
        output artefacts are copied from the container's temporary directory to a staging area on
        the \gls{vm} filesystem via \texttt{docker cp}.
  \item Result download (\gls{sftp} GET): the \gls{sftp} client iterates over the
        staging directory and downloads each file to the local \texttt{outputs/} directory.
\end{enumerate}

This design avoids mounting any host directories into the Docker containers on the \gls{vm}:
all file movements cross container boundaries via \texttt{docker cp}, and all client--\gls{vm}
movements use \gls{sftp}.  The \gls{csv} files themselves are never transferred from the
\gls{vm} to the client; only the small result artefacts (aggregated \gls{csv}, scatter plot,
\gls{json} summary) are downloaded.

% ─────────────────────────────────────────────────────────────────────────────
\section{End-to-End Data Flow}
\label{sec:rm-dataflow}
% ─────────────────────────────────────────────────────────────────────────────

Figure~\ref{fig:rm-dataflow} presents the complete end-to-end data flow of the remote deployment
from the analyst's perspective.

\begin{figure}[htbp]
\centering
\scalebox{0.85}{%
\begin{tikzpicture}[x=1cm, y=1cm,
  anabox/.style={rectangle, rounded corners=3pt, draw=black!70, fill=cyan!8,
                 text width=4.0cm, minimum height=0.9cm, text centered,
                 font=\small, inner sep=5pt},
  vmbox/.style={rectangle, rounded corners=3pt, draw=black!70, fill=orange!10,
                text width=4.0cm, minimum height=0.9cm, text centered,
                font=\small, inner sep=5pt},
  cntbox/.style={rectangle, rounded corners=3pt, draw=black!70, fill=yellow!10,
                 text width=3.5cm, minimum height=0.9cm, text centered,
                 font=\small, inner sep=5pt},
  fsbox/.style={rectangle, rounded corners=3pt, draw=black!60, fill=gray!8,
                text width=3.5cm, minimum height=0.9cm, text centered,
                font=\small, inner sep=5pt},
  carr/.style={thick,->,>=Stealth},
  sarr/.style={thick,->,>=Stealth,dashed,blue!70},
  farr/.style={thick,->,>=Stealth,dotted,green!60!black},
]

%% Column headers
\node[font=\bfseries\small,cyan!60!black] at (-5, 2.0) {Analyst Machine};
\node[font=\bfseries\small,orange!70!black] at (1.8, 2.0) {VM Filesystem};
\node[font=\bfseries\small,yellow!50!black] at (7.5, 2.0) {Docker Containers};

%% Analyst side
\node[anabox] (yaml)   at (-5,  0.5) {Local YAML scenario};
\node[anabox] (client) at (-5, -1.0) {cli.py (paramiko)};
\node[fsbox]  (out)    at (-5, -3.0) {outputs/ (results)};

%% VM side
\node[vmbox]  (vmtmp)  at (1.8,  0.5) {/tmp/ (YAML staging)};
\node[vmbox]  (vmdl)   at (1.8, -3.0) {/tmp/aerosouda\_dl/};

%% Container side
\node[cntbox] (master) at (7.5,  0.5) {spark-master};
\node[cntbox] (hdfs)   at (7.5, -1.0) {HDFS (namenode + datanodes)};
\node[cntbox] (workers)at (7.5, -2.7) {spark-worker-1..6};
\node[cntbox] (cntout) at (7.5, -4.2) {spark-master:\\{}/tmp/outputs/};

%% Arrows
\draw[sarr]  (yaml.east)   -- node[above,font=\scriptsize]{SFTP PUT} (vmtmp.west);
\draw[sarr]  (vmtmp.east)  -- node[above,font=\scriptsize]{docker exec} (master.west);
\draw[sarr]  (client.east) -- node[below,font=\scriptsize]{SSH: list /inputs/} (hdfs.west);
\draw[sarr]  (client.east) -- node[below,font=\scriptsize]{SSH: spark-submit} (master.west);
\draw[carr]  (hdfs.south)  -- node[left,font=\scriptsize,pos=0.5,xshift=2pt]{HDFS read} (workers.north);
\draw[sarr]  (client.east) -- node[below,font=\scriptsize,fill=white,inner sep=1pt,pos=0.4]{docker cp} (vmdl.west);
\draw[carr]  (cntout.west) -- (vmdl.east);
\draw[farr]  (vmdl.west)   -- node[below,font=\scriptsize]{SFTP GET} (out.east);

\end{tikzpicture}%
}
\caption[End-to-end data flow of the AeroSOuDA remote deployment.]{End-to-end data flow of the AeroSOuDA remote deployment.
         Dashed blue arrows indicate SSH-mediated commands; dotted green arrows indicate
         SFTP file transfers; solid black arrows indicate direct data reads within
         the Docker network.}
\label{fig:rm-dataflow}
\end{figure}

The key insight of the design is the asymmetry between the upward and downward data flows.
In the upward direction (analyst $\to$ \gls{vm}), only the small \gls{yaml} scenario file
(typically a few kilobytes) is transferred.  The large simulation \gls{csv} files are already
resident on the \gls{vm}'s \gls{hdfs} storage layer, uploaded once by a \gls{vm} operator.
In the downward direction (\gls{vm} $\to$ analyst), only the compact result artefacts are
downloaded: an aggregated \gls{csv} (at most a few hundred rows), a scatter plot image, and a
\gls{json} summary.  This design minimises network bandwidth consumption and makes the pipeline
practical even on a slow connection.

The simulation \gls{csv} data never leaves the \gls{vm}.  It is ingested by the Spark workers
directly from \gls{hdfs} within the Docker bridge network, processed in memory across the six
worker containers, and the aggregated results are written to the master's temporary directory.
Only the aggregated output is exfiltrated to the analyst's machine.

\subsection{Separation of Concerns}
\label{subsec:rm-separation-of-concerns}

The two-repository design enforces a clean separation of concerns:

\begin{itemize}
  \item Cluster (\texttt{tg-cluster/cluster/}): responsible for storage and
        computation.  It manages the \gls{hdfs} data layer and the Spark execution engine.  Its
        only external interface is a standard Docker Compose stack accessible via \gls{ssh}.
  \item Client (\texttt{tg-cluster/client/}): responsible for user interaction
        and orchestration.  It knows the analyst's intent (which \gls{csv}, which scenario) and
        coordinates all actions on the \gls{vm} cluster via \gls{ssh} commands.  It has no
        knowledge of Spark internals or \gls{hdfs} block placement.
\end{itemize}

This separation means that the cluster can be upgraded (more workers, larger DataNodes, different
Spark version) without any changes to the client, and the client can be updated (new scenario
wizard features, different output formats) without touching the cluster configuration.

\subsection{Execution Phase Instrumentation}
\label{subsec:rm-phase-instrumentation}

The \texttt{scripts/processor/spark\_job.py} module emits fine-grained timing markers at
key boundaries of the pipeline.  Each marker is printed to \texttt{stdout} with the prefix
\texttt{METRIC:}, for example:

\begin{verbatim}
METRIC:T_INIT_SECONDS=18.432160
METRIC:T_MAP_SECONDS=55.117832
METRIC:T_REDUCE_SECONDS=0.009412
METRIC:T_OUTPUT_SECONDS=0.421073
\end{verbatim}

The four markers decompose $T_{\text{job}}$ into four sequential phases:

\begin{itemize}
  \item $T_{\text{init}}$: time from \texttt{SparkSession.getOrCreate()} returning
        to the first data-read call.  Captures \gls{jvm} initialisation, Spark context
        setup, and worker registration.
  \item $T_{\text{map}}$: time for the \texttt{mapPartitions().collect()} action.
        This is the dominant phase and encompasses both \gls{hdfs} data reads (I/O) and
        distributed Python payload parsing (computation).  Because Spark evaluates plans
        lazily, these two sub-phases cannot be separated without an additional materialise
        step (e.g., \texttt{cache().count()}) that would change the benchmark semantics.
  \item $T_{\text{reduce}}$: time for the driver-side reduce loop that merges the
        per-partition partial results into the global accumulators.  This phase is
        $O(n_{\text{entities}})$ and is expected to be negligible.
  \item $T_{\text{output}}$: time for constructing the output Pandas structure,
        writing the aggregated \gls{csv}, rendering the scatter plot, and serialising the
        \gls{json} summary.
\end{itemize}

The \texttt{tg-cluster} benchmark harness parses these markers in real time as the
\gls{ssh} output stream arrives and records them as additional columns in the benchmark
\gls{csv} (\texttt{t\_init\_s}, \texttt{t\_map\_s}, \texttt{t\_reduce\_s},
\texttt{t\_output\_s}).  In parallel, a background sampling thread issues periodic
\texttt{docker stats} queries over the same \gls{ssh} connection to record peak and mean
worker \gls{cpu} usage (\texttt{cpu\_peak\_pct}, \texttt{cpu\_mean\_pct}) and peak
memory consumption (\texttt{mem\_peak\_mb}) across the four Spark worker containers.

This instrumentation provides the empirical basis for quantifying host-resource contention
at scale and provides enough content to test and analyse this thesis' capacities in the next
session.
